% GENERATED FILE. Edit canonical lessons, then run npm run generate:books.
% canonical-source-hash: fnv1a64:e7ed5b3af0ecf864
% canonical-lessons: TE-C498-atmagauravam, TE-C498-avamanam, TE-C498-garvapadu, TE-C498-mohamatapadu, TE-C498-taladincuko

\chapter{Shame and Shyness}
\label{ch:te-shame-and-shyness}
\hlchaptermodality{498}

\begin{chapteropening}
\textbf{By the end of this chapter:} \emph{I can talk about self-respect, humiliation and pride, holding back and hanging my head.}

Complete the last lesson of chapter 498: I can talk about self-respect, humiliation and pride, holding back and hanging my head.
\end{chapteropening}

\section[\texorpdfstring{\te{ఆత్మగౌరవం} (ātmagauravaṁ)}{ఆత్మగౌరవం (ātmagauravaṁ)}]{\te{ఆత్మగౌరవం} (ātmagauravaṁ) — self-respect; dignity}

\label{lesson:TE-C498-atmagauravam}



\begin{quote}
Before the new one: say the Telugu for fury, then the Telugu for irritation.
\end{quote}

\subsection*{You'll want to know: \te{ఆత్మగౌరవం}}
\textbf{\te{ఆత్మగౌరవం}} — \emph{ātmagauravaṁ} — \textquotedblleft{}self-respect; dignity\textquotedblright{}.

\begin{quote}
\te{పేదవాడైనా} \te{అతనికి} \te{ఆత్మగౌరవం} \te{ఎక్కువ}.
\end{quote}

\begin{quote}
\emph{pēdavāḍainā ataniki ātmagauravaṁ ekkuva.}
\end{quote}

\begin{quote}
Though he is poor, he has a strong sense of self-respect.
\end{quote}

\te{ఆత్మ}, self, + \te{గౌరవం}, respect. \te{ఆత్మగౌరవంతో} \te{బతుకు}, to live with dignity.

\begin{cousinweb}
\textbf{\te{ఆత్మగౌరవం}} (\emph{ātmagauravaṁ}) is a learned compound built from Sanskrit \emph{ātman}, \textquotedblleft{}breath, self\textquotedblright{} (CDIAL 1135), and \emph{gaurava}, \textquotedblleft{}weight; venerableness\textquotedblright{} (CDIAL 4346), the word behind \te{గౌరవం}, \textquotedblleft{}respect\textquotedblright{}.
\end{cousinweb}

\begin{grammarlens}[title={What you've built}]
One more word for describing what you lived and felt.
\end{grammarlens}

\subsection*{Your turn}
\begin{itemize}
\raggedright
  \item \emph{Say it:} \emph{ātmagauravaṁ}
  \item \emph{Say it:} \emph{ātmagauravaṁ}, once more
  \item \emph{Say it:} \emph{pēdavāḍainā ataniki ātmagauravaṁ ekkuva.} — Though he is poor, he has a strong sense of self-respect.
  \item \emph{Recall:} say the Telugu for fury, then the Telugu for irritation, then say \emph{ātmagauravaṁ} again
\end{itemize}

\begin{tcolorbox}[breakable,colback=teal!4,colframe=teal!35!black,title={Before you move on}]
What does \te{ఆత్మగౌరవం} mean? (\textquotedblleft{}Self-respect; dignity\textquotedblright{}.) Say it once more.
\end{tcolorbox}
\section[\texorpdfstring{\te{అవమానం} (avamānaṁ)}{అవమానం (avamānaṁ)}]{\te{అవమానం} (avamānaṁ) — humiliation; an insult}

\label{lesson:TE-C498-avamanam}



\begin{quote}
Before the new one: say the Telugu for irritation, then the Telugu for self-respect.
\end{quote}

\subsection*{You'll want to know: \te{అవమానం}}
\textbf{\te{అవమానం}} — \emph{avamānaṁ} — \textquotedblleft{}humiliation; an insult\textquotedblright{}.

\begin{quote}
\te{అందరి} \te{ముందు} \te{తిట్టడం} \te{అతనికి} \te{అవమానంగా} \te{అనిపించింది}.
\end{quote}

\begin{quote}
\emph{andari mundu tiṭṭaḍaṁ ataniki avamānaṅgā anipiñcindi.}
\end{quote}

\begin{quote}
Being scolded in front of everyone felt humiliating to him.
\end{quote}

Its opposite is \te{గౌరవం}, respect.

\begin{cousinweb}
\textbf{\te{అవమానం}} (\emph{avamānaṁ}) is borrowed from Sanskrit \emph{avamāna}, \textquotedblleft{}disrespect, dishonour\textquotedblright{} (Monier-Williams), from \emph{ava-}, \textquotedblleft{}down\textquotedblright{}, and \emph{man}, \textquotedblleft{}to think\textquotedblright{}: to think someone low.
\end{cousinweb}

\begin{grammarlens}[title={What you've built}]
One more word for describing what you lived and felt.
\end{grammarlens}

\subsection*{Your turn}
\begin{itemize}
\raggedright
  \item \emph{Say it:} \emph{avamānaṁ}
  \item \emph{Say it:} \emph{avamānaṁ}, once more
  \item \emph{Say it:} \emph{andari mundu tiṭṭaḍaṁ ataniki avamānaṅgā anipiñcindi.} — Being scolded in front of everyone felt humiliating to him.
  \item \emph{Recall:} say the Telugu for irritation, then the Telugu for self-respect, then say \emph{avamānaṁ} again
\end{itemize}

\begin{tcolorbox}[breakable,colback=teal!4,colframe=teal!35!black,title={Before you move on}]
What does \te{అవమానం} mean? (\textquotedblleft{}Humiliation; an insult\textquotedblright{}.) Say it once more.
\end{tcolorbox}
\section[\texorpdfstring{\te{గర్వపడు} (garvapaḍu)}{గర్వపడు (garvapaḍu)}]{\te{గర్వపడు} (garvapaḍu) — to be proud}

\label{lesson:TE-C498-garvapadu}



\begin{quote}
Before the new one: say the Telugu for self-respect, then the Telugu for humiliation.
\end{quote}

\subsection*{You'll want to know: \te{గర్వపడు}}
\textbf{\te{గర్వపడు}} — \emph{garvapaḍu} — \textquotedblleft{}to be proud\textquotedblright{}.

\begin{quote}
\te{కూతురి} \te{విజయం} \te{చూసి} \te{తల్లిదండ్రులు} \te{గర్వపడ్డారు}.
\end{quote}

\begin{quote}
\emph{kūturi vijayaṁ cūsi tallidaṇḍrulu garvapaḍḍāru.}
\end{quote}

\begin{quote}
Seeing their daughter succeed, the parents were proud.
\end{quote}

Past (I …): \te{గర్వపడ్డాను}, \emph{garvapaḍḍānu}.

The noun \te{గర్వం} can also mean arrogance; the verb is the good kind of pride.

\begin{cousinweb}
\textbf{\te{గర్వపడు}} (\emph{garvapaḍu}) is built from \te{గర్వం}, \textquotedblleft{}pride\textquotedblright{}, from Sanskrit \emph{garva}, \textquotedblleft{}pride, arrogance\textquotedblright{} (CDIAL 4064) and the native verb \emph{paḍu}, \textquotedblleft{}to feel, to suffer\textquotedblright{} (DEDR 3853).
\end{cousinweb}

\begin{grammarlens}[title={What you've built}]
One more word for describing what you lived and felt.
\end{grammarlens}

\subsection*{Your turn}
\begin{itemize}
\raggedright
  \item \emph{Say it:} \emph{garvapaḍu}
  \item \emph{Say it:} \emph{garvapaḍu}, once more
  \item \emph{Say it:} \emph{kūturi vijayaṁ cūsi tallidaṇḍrulu garvapaḍḍāru.} — Seeing their daughter succeed, the parents were proud.
  \item \emph{Recall:} say the Telugu for self-respect, then the Telugu for humiliation, then say \emph{garvapaḍu} again
\end{itemize}

\begin{tcolorbox}[breakable,colback=teal!4,colframe=teal!35!black,title={Before you move on}]
What does \te{గర్వపడు} mean? (\textquotedblleft{}To be proud\textquotedblright{}.) Say it once more.
\end{tcolorbox}
\section[\texorpdfstring{\te{మొహమాటపడు} (mohamāṭapaḍu)}{మొహమాటపడు (mohamāṭapaḍu)}]{\te{మొహమాటపడు} (mohamāṭapaḍu) — to hold back politely; to be too polite to ask}

\label{lesson:TE-C498-mohamatapadu}



\begin{quote}
Before the new one: say the Telugu for humiliation, then the Telugu for to be proud.
\end{quote}

\subsection*{You'll want to know: \te{మొహమాటపడు}}
\textbf{\te{మొహమాటపడు}} — \emph{mohamāṭapaḍu} — \textquotedblleft{}to hold back politely; to be too polite to ask\textquotedblright{}.

\begin{quote}
\te{అతను} \te{మొహమాటపడి} \te{ఇంకా} \te{అన్నం} \te{అడగలేదు}.
\end{quote}

\begin{quote}
\emph{atanu mohamāṭapaḍi iṅkā annaṁ aḍagalēdu.}
\end{quote}

\begin{quote}
Out of politeness he did not ask for more rice.
\end{quote}

Past (I …): \te{మొహమాటపడ్డాను}, \emph{mohamāṭapaḍḍānu}.

The noun \te{మొహమాటం} is the polite reluctance to refuse, to ask or to say no.

\begin{cousinweb}
\textbf{\te{మొహమాటపడు}} (\emph{mohamāṭapaḍu}) is built from \te{మొహమాటం}, the polite reluctance to refuse or to ask, and the native verb \emph{paḍu}, \textquotedblleft{}to feel, to suffer\textquotedblright{} (DEDR 3853). The origin of \te{మొహమాటం} is not settled.
\end{cousinweb}

\begin{grammarlens}[title={What you've built}]
One more word for describing what you lived and felt.
\end{grammarlens}

\subsection*{Your turn}
\begin{itemize}
\raggedright
  \item \emph{Say it:} \emph{mohamāṭapaḍu}
  \item \emph{Say it:} \emph{mohamāṭapaḍu}, once more
  \item \emph{Say it:} \emph{atanu mohamāṭapaḍi iṅkā annaṁ aḍagalēdu.} — Out of politeness he did not ask for more rice.
  \item \emph{Recall:} say the Telugu for humiliation, then the Telugu for to be proud, then say \emph{mohamāṭapaḍu} again
\end{itemize}

\begin{tcolorbox}[breakable,colback=teal!4,colframe=teal!35!black,title={Before you move on}]
What does \te{మొహమాటపడు} mean? (\textquotedblleft{}To hold back politely; to be too polite to ask\textquotedblright{}.) Say it once more.
\end{tcolorbox}
\section[\texorpdfstring{\te{తలదించుకో} (taladiñcukō)}{తలదించుకో (taladiñcukō)}]{\te{తలదించుకో} (taladiñcukō) — to hang one's head; to be ashamed}

\label{lesson:TE-C498-taladincuko}



\begin{quote}
Before the new one: say the Telugu for to be proud, then the Telugu for to hold back politely.
\end{quote}

\subsection*{You'll want to know: \te{తలదించుకో}}
\textbf{\te{తలదించుకో}} — \emph{taladiñcukō} — \textquotedblleft{}to hang one's head; to be ashamed\textquotedblright{}.

\begin{quote}
\te{అబద్ధం} \te{బయటపడగానే} \te{అతను} \te{తలదించుకున్నాడు}.
\end{quote}

\begin{quote}
\emph{abaddhaṁ bayaṭapaḍagānē atanu taladiñcukunnāḍu.}
\end{quote}

\begin{quote}
As soon as the lie came out, he hung his head.
\end{quote}

Past (I …): \te{తలదించుకున్నాను}, \emph{taladiñcukunnānu}.

\te{తల}, head, + \te{దించుకో}, to lower your own: the gesture of shame. Its opposite is \te{తలెత్తుకో}, to hold your head high.

\begin{cousinweb}
\textbf{\te{తలదించుకో}} (\emph{taladiñcukō}) is native Dravidian: \te{తల}, \textquotedblleft{}head\textquotedblright{} (DEDR 3103), and \te{దించు}, \textquotedblleft{}to lower, to put down\textquotedblright{} (DEDR 502), with -\te{కో}, \textquotedblleft{}for oneself\textquotedblright{}, from \te{కొను}, \textquotedblleft{}to take\textquotedblright{} (DEDR 2151): to lower one's own head.
\end{cousinweb}

\begin{grammarlens}[title={What you've built}]
One more word for describing what you lived and felt.
\end{grammarlens}

\subsection*{Your turn}
\begin{itemize}
\raggedright
  \item \emph{Say it:} \emph{taladiñcukō}
  \item \emph{Say it:} \emph{taladiñcukō}, once more
  \item \emph{Say it:} \emph{abaddhaṁ bayaṭapaḍagānē atanu taladiñcukunnāḍu.} — As soon as the lie came out, he hung his head.
  \item \emph{Recall:} say the Telugu for to be proud, then the Telugu for to hold back politely, then say \emph{taladiñcukō} again
\end{itemize}

\begin{tcolorbox}[breakable,colback=teal!4,colframe=teal!35!black,title={Before you move on}]
What does \te{తలదించుకో} mean? (\textquotedblleft{}To hang one's head; to be ashamed\textquotedblright{}.) Say it once more.
\end{tcolorbox}
